% ── §3  Motivation ────────────────────────────────────────────────────────────
\section{Motivation: Where Does Tail Latency Come From?}
\label{sec:motivation}

To motivate \xtal, we characterise the uncontrolled Kafka\,+\,Flink
pipeline under three load patterns and four static configurations.
This study answers two questions:
(1)~Under what conditions does tail latency spike, and how quickly?
(2)~Which internal signal (\texttt{inPoolUsage} vs.\
checkpoint duration vs.\ consumer lag) spikes first?

\subsection{Experimental Setup}
We run Flink 1.18 (1 JobManager, 3 TaskManagers, 2 slots each) and
Kafka 3.6 in Docker on a single host.
The workload is NEXMark Query~5 (sliding window auction join) driven
by our load generator at 5{,}000~events/s.
Metrics are scraped via Prometheus at 5-second intervals.
We test four static configurations:
\begin{itemize}
  \item \textbf{Default:} checkpoint=5\,s, buffer\_timeout=100\ms, linger=5\ms
  \item \textbf{Low-latency:} checkpoint=500\ms, buffer\_timeout=0\ms, linger=0\ms
  \item \textbf{High-throughput:} checkpoint=10\,s, buffer\_timeout=100\ms, linger=50\ms
  \item \textbf{Hand-tuned:} checkpoint=2\,s, buffer\_timeout=10\ms, linger=5\ms
\end{itemize}

\subsection{Steady Load}
Under constant 5{,}000~events/s, all configurations achieve \pnn{99}
below the 500\ms SLA, with the default configuration exhibiting
\pnn{99}\,$\approx$\,180\ms and periodic spikes to 350\ms coinciding
with checkpoint barriers (visible as \texttt{lastCheckpointDuration}
pulses in Prometheus).

\subsection{Bursty Load}
At $t=30$\,s, throughput jumps 10$\times$ to 50{,}000~events/s for 10\,s.
The sequence of events is consistent across all configurations:
\begin{enumerate}
  \item Kafka consumer lag grows within 1\,s as Flink cannot drain the
    partition at burst rate.
  \item Input buffer usage (\texttt{inPoolUsage}) saturates to $>$0.9
    within 5\,s, triggering Flink backpressure.
  \item Checkpoint barriers stall behind full buffers;
    \texttt{lastCheckpointDuration} spikes from $\sim$850\ms to
    $>$8{,}000\ms.
  \item End-to-end \pnn{99} rises to 1{,}200--3{,}400\ms (2.4--6.8$\times$
    the SLA) and remains elevated for 25--60\,s after the burst ends.
\end{enumerate}
The \textbf{hand-tuned} configuration recovers fastest (28\,s) because its
shorter checkpoint interval (2\,s) limits barrier alignment overhead.
The \textbf{high-throughput} configuration is worst, taking 62\,s to recover.

\textbf{Finding~1:} Consumer lag and input buffer usage are leading
indicators of a tail spike; checkpoint duration spikes are the mechanism
that sustains the spike.

\subsection{Failure Injection}
We kill one TaskManager at $t=45$\,s.
Flink detects the failure and triggers job restart with a savepoint.
During recovery (30--50\,s), \pnn{99} is undefined (no results produced).
After recovery, the two remaining TaskManagers are overloaded:
consumer lag grows and \pnn{99} settles at 800--1{,}400\ms---above SLA---
until the operator is restarted with reduced parallelism or the burst abates.

\textbf{Finding~2:} Post-failure, the pipeline often stabilises at a new
operating point \emph{above} the SLA without intervention. A controller
that can reduce checkpoint frequency and buffer timeout post-recovery
accelerates the return to SLA compliance.
